\clearpage\SourceLeaf{137}
\begin{center}
[III.]

\medskip
\textbf{MUSTERRECHNUNG, UM AUS $A=p\cos P$, $B=p\sin P$\\
$p$ UND $P$ ZU FINDEN.}
\end{center}

\begin{center}\small
\begin{tabular}{l@{\;}l@{\qquad\hspace{.8em}}l}
$B\; .\ .\ .\ .\ .\ .$ & $9{,}56905\;69225\,n$ &\\
$A\; .\ .\ .\ .\ .\ .$ & $0{,}28523\;04177\,n$ &\\\cmidrule(r{1.4em}){1-2}
$\mathrm{tg}\,P\; .\ .\ .\ .\ .$ & $9{,}28382\;65048$ & $P=190^0\,52'\,52''{,}913567$\\
$\mathrm{tg}\,P^*\; .\ .\ .\ .$ & $9{,}28379\;34121$ & $P^*=190\;52\;50$\\\cmidrule(r{1.4em}){1-2}
$\cos P^*\; .\ .\ .$ & $9{,}99212\;15841\,n$ &\\
& $-11792$ & $5{,}5197322\; .\ .\ .\ .\ .\;3\;30927$\\
$\cos P\; .\ .\ .\ .$ & $9{,}99212\;04049\,n$ & $\underline{9{,}9842419}\; .\ .\ .\ .\;\cos P\cos P^*$\\\cmidrule(r{1.4em}){1-2}
$p\; .\ .\ .\ .\ .\ .\ .$ & $0{,}29311\;00128$ & $5{,}5039741$
\end{tabular}
\end{center}

\begin{center}\small
\renewcommand{\arraystretch}{1.35}
\begin{tabular}{r@{\;}l@{\;}l}
\multirow{2}{*}{$11792\Bigg\{$}&$9{,}2838100\; .\ .\ .\ .\;\BareRoot{(\mathrm{tg}\,P\mathrm{tg}\,P^*)}$&\\
&$\overline{4{,}7877841}$&\multirow{2}{*}{$\Bigg\}\;2''{,}913567.$}\\
$\left[-\log\frac{\mathrm{Mod.}}{206264{,}8}\right]=\mathrm{Const.}\; .\ .\ .$&$-4{,}3233592$&
\end{tabular}
\end{center}

\medskip
\begin{center}\textbf{BEMERKUNGEN.}\end{center}

Dass die Anzeige, S. 121--127, von \textsc{Gauss} herrührt, geht aus einem Briefe von
\textsc{Gauss} an \textsc{Olbers} vom 3. December 1808 hervor, worin es heisst:

\begin{quotation}\itshape
„So wünschte ich z. B. sehr, dass eine solche Tafel, wie \textup{\textsc{Leonelli}}
vorgeschlagen hat (meine Anzeige: Hallische L.Z. 1808 vom 12. Febr.), ausgeführt würde.
Für bloss 5 Decimalen habe ich selbst einmal einen Anfang gemacht. Bei solchen
Rechnungen, wo sehr viele Logarithmen von Summen oder Differenzen gesucht
werden (wie z. B. bei meiner Methode die Störungen zu berechnen), würde eine
solche Tafel eine bedeutende Erleichterung geben«.
\end{quotation}

Die Notiz [I.] ist einem Handbuch entnommen; [II.] und [III.] finden sich auf Blättern,
die den Büchern der \textsc{Gauss}-Bibliothek: \textsc{Hobert} und \textsc{Ideler},
Trigonometrische Tafeln, und \textsc{Vega}, Thesaurus Logarithmorum, angeheftet sind.
\hfill \textsc{Börsch, Krüger}.
